% Part: first-order-logic
% Chapter: models-theories
% Section: size-of-structures

\documentclass[../../../include/open-logic-section]{subfiles}

\begin{document}

\olfileid[fr]{fol}{mat}{siz}

\olsection{Exprimer la taille des structures}

\begin{explain}
Certaines propriétés des structures s'expriment sans même recourir
aux symboles non logiques du langage. Il existe, par exemple, des
!!{sentence}s qui sont vrais dans !!a{structure} si et seulement si
le !!{domain} de cette !!{structure} possède au moins, au plus ou
exactement un certain nombre~$n$ d'!!{element}s.\footnote{Précision éditoriale : les affichages avec des points de suspension indiquent les comparaisons entre toutes les paires d'indices distincts. Pour un seul indice, la conjonction sans comparaison vaut vrai : on peut prendre $!A_{\ge 1}\ident\lexists[x][\eq[x][x]]$ et $!A_{=1}\ident\lexists[x][\lforall[y][\eq[y][x]]]$. Pour zéro, $!A_{\ge0}$ est toujours vrai et $!A_{=0}$ toujours faux, puisque les domaines sont non vides.}
\end{explain}

\begin{prop}
L'!!{sentence}
\begin{multline*}
  !A_{\ge n} \ident \lexists[x_1][\lexists[x_2][\dots\lexists[x_n][{}]]]\\
  \begin{aligned}
  (\eq/[x_1][x_2] \land {}
  \eq/[x_1][x_3] \land \eq/[x_1][x_4] \land \dots \land \eq/[x_1][x_n] \land {}\\
\eq/[x_2][x_3] \land \eq/[x_2][x_4] \land \dots \land {} \eq/[x_2][x_n] \land {} \\
\vdots\\
\eq/[x_{n-1}][x_n])
\end{aligned}
\end{multline*}
est vrai dans !!a{structure}~$\Struct M$ si et seulement si
$\Domain M$ contient au moins $n$~!!{element}s. Par conséquent,
$\Sat{M}{\lnot !A_{\ge n+1}}$ si et seulement si $\Domain M$
contient au plus~$n$~!!{element}s.

\end{prop}

\begin{prop}
L'!!{sentence}
\begin{multline*}
!A_{= n} \ident \lexists[x_1][\lexists[x_2][\dots\lexists[x_n][{}]]] \\
\begin{aligned}
  (\eq/[x_1][x_2] \land {}
  \eq/[x_1][x_3] \land \eq/[x_1][x_4] \land \dots \land \eq/[x_1][x_n] \land {}\\
\eq/[x_2][x_3] \land \eq/[x_2][x_4] \land \dots \land {} \eq/[x_2][x_n] \land {} \\
\vdots\\
\eq/[x_{n-1}][x_n] \land {} \\
\lforall[y][(\eq[y][x_1] \lor \dots \lor \eq[y][x_n]]))
\end{aligned}
\end{multline*}
est vrai dans !!a{structure}~$\Struct M$ si et seulement si
$\Domain M$ contient exactement $n$~!!{element}s.
\end{prop}

\begin{prop}
Une !!{structure} est infinie si et seulement si elle est un modèle
de
\[
\{!A_{\ge 1}, !A_{\ge 2}, !A_{\ge 3}, \dots \}.
\]
\end{prop}

Aucun énoncé purement logique ne peut être vrai dans~$\Struct M$ si
et seulement si $\Domain M$ est infini. En revanche, il existe des
!!{sentence}s comportant des !!{predicate}s non logiques qui
n'admettent que des modèles infinis, sans que toute !!{structure}
infinie en soit nécessairement un modèle. La propriété d'être une
structure finie et celle d'être une structure !!{nonenumerable} ne
peuvent pas même être exprimées par un ensemble infini
d'!!{sentence}s. Ces faits découlent des théorèmes de compacité et
de Löwenheim--Skolem.

\end{document}
